% !TeX spellcheck = ru_RU
\section{Постановка задачи}
\label{sec:task}

Целью работы является разработка системы извлечения и анализа связей между персонажами и локациями из текстов фандомных вики-источников с обеспечением гибкой архитектуры, позволяющей заменять источник данных без изменения аналитической части.

Для достижения поставленной цели были решены следующие задачи:
\begin{enumerate}
    \item Спроектировать модульную архитектуру, в которой бэкенд (парсинг) и DS-часть (извлечение сущностей, аналитика) не зависят друг от друга, а парсер конкретного фандома можно заменить, написав новый класс, реализующий единый интерфейс.
    \item Реализовать парсер для выбранного фандома (на примере вселенной «Властелин колец»), который собирает тексты страниц, очищает их от HTML-тегов и сохраняет в CSV-файл.
    \item Разработать модуль извлечения именованных сущностей (NER), выделяющий персонажей (тип PERSON) и локации (тип LOC) из текстов, присваивающий каждой уникальной сущности идентификатор и формирующий таблицы связей «документ–сущность».
    \item Загрузить полученные CSV-файлы в реляционную базу данных PostgreSQL, обеспечив корректную обработку автоинкрементных полей и внешних ключей.
    \item Реализовать SQL-запросы для получения:
        \begin{itemize}
            \item топ-10 персонажей заданной локации;
            \item списка локаций заданного персонажа;
            \item персонажей, часто упоминаемых вместе с заданным;
            \item списка документов, содержащих заданную сущность.
        \end{itemize}
    \item Проверить адекватность системы на реальных данных (несколько страниц вики) и продемонстрировать работу запросов.
\end{enumerate}

Результатом работы является программный комплекс (исходный код на Python, SQL-скрипты) и описание архитектуры, позволяющее легко адаптировать систему к любому другому фандому.
